\clearpage
\SourceLeaf{411}%
\fancyhead[L]{\small MISCELLANEOUS PAPERS.}
\pdfbookmark[1]{Solution of a geometrical problem.}{aufloesung-geometrische-1823}
\begin{center}
\rule{.8\textwidth}{.4pt}\\[4pt]
Astronomische Nachrichten. No. 42. November 1823.\\[4pt]
\rule{.8\textwidth}{.4pt}\\[12pt]
\textit{Solution of a geometrical problem.}
\end{center}
\hspace*{\parindent}\{On p.\,61 of Professor \textsc{Möbius}'s description of the Leipzig Observatory there is a small geometrical problem, as follows.
\begin{quote}
Take any 5 points \(A,B,C,D,E\) in a plane and join each pair by a straight line. The areas of the 5 resulting triangles \(EAB,ABC,BCD,CDE,DEA\) are known, and we are to find from them the area of the pentagon \(ABCDE\).
\end{quote}
Court Councillor \textsc{Gauss}, who spent a few weeks with me this summer, wrote the following solution in it when he saw the book.

Denote the 5 points by \(1,2\ 3,4,5\), the
\begin{center}
\begin{tabular}{lrl}
angles&\(213\)&by \(p\)\\
&\(214\)&— \(q\)\\
&\(215\)&— \(r\)\\\hline
sides&\(12\)&— \(t\)\\
&\(13\)&— \(u\)\\
&\(14\)&— \(v\)\\
&\(15\)&— \(w\)\\\hline
triangles&\(123\)&— \(a\)\\
&\(234\)&— \(b\)\\
&\(345\)&— \(c\)\\
&\(451\)&— \(d\)\\
&\(512\)&— \(e\)\\
&\(124\)&— \(x\)\\
&\(134\)&— \(y\)\\
&\(135\)&— \(z\)\\\hline
pentagon&\(12345\)&— \(\omega\)
\end{tabular}
\end{center}
